\section*{Розділ VIII. Фінансування та майно}
\addcontentsline{toc}{section}{Розділ VIII. Фінансування та майно}

\subsection*{8.1. Джерела фінансування ЦВКс}
\addcontentsline{toc}{subsection}{8.1. Джерела фінансування ЦВКс}
    8.1.1. Фінансування діяльності ЦВКс здійснюється за рахунок коштів, передбачених у єдиному кошторисі органів студентського самоврядування Університету відповідно до Розділу V Положення про ОСС.

    8.1.2. ЦВКс має право залучати додаткові кошти для забезпечення виборчого процесу у формі:

        \begin{enumerate}[label=\alph*)]
            \item Цільових внесків від органів студентського самоврядування;
            \item Технічної підтримки від адміністрації Університету;
            \item Благодійної допомоги відповідно до законодавства України.
        \end{enumerate}

    8.1.3. Залучення додаткових коштів здійснюється прозоро та не повинно впливати на незалежність ЦВКс у прийнятті рішень.

\subsection*{8.2. Планування та використання коштів}
\addcontentsline{toc}{subsection}{8.2. Планування та використання коштів}
    8.2.1. ЦВКс щорічно подає до КСУ кошторис витрат на забезпечення виборчих процесів та поточної діяльності.

    8.2.2. Кошти ЦВКс використовуються виключно на цілі, передбачені цим Положенням, зокрема:

        \begin{enumerate}[label=\alph*)]
            \item Організаційно-технічне забезпечення виборів;
            \item Виготовлення виборчих матеріалів та документів;
            \item Оренда приміщень для проведення виборів (у разі необхідності);
            \item Технічне обслуговування електронних систем голосування;
            \item Інформаційне забезпечення виборчого процесу;
            \item Підвищення кваліфікації членів ЦВКс.
        \end{enumerate}

    8.2.3. Усі фінансові операції ЦВКс здійснюються через бухгалтерію Університету відповідно до встановлених процедур фінансової діяльності.

\subsection*{8.3. Контроль за використанням коштів}
\addcontentsline{toc}{subsection}{8.3. Контроль за використанням коштів}
    8.3.1. ЦВКс веде облік всіх витрат, пов'язаних з виборчим процесом, та готує фінансові звіти за кожними виборами.

    8.3.2. Фінансову звітність ЦВКс перевіряє СУД у рамках своїх контрольних повноважень.

    8.3.3. Річний фінансовий звіт ЦВКс подається до КСУ разом із звітом про діяльність.

    8.3.4. Інформація про витрати на проведення виборів оприлюднюється на офіційних інформаційних ресурсах органів студентського самоврядування.

\subsection*{8.4. Майнове забезпечення}
\addcontentsline{toc}{subsection}{8.4. Майнове забезпечення}
    8.4.1. Адміністрація Університету надає ЦВКс необхідні приміщення та технічні засоби для здійснення її діяльності відповідно до Положення про ОСС.

    8.4.2. ЦВКс має право користуватися:

        \begin{enumerate}[label=\alph*)]
            \item Приміщеннями для проведення засідань та виборів;
            \item Офісною технікою та засобами зв'язку;
            \item Електронними системами та програмним забезпеченням;
            \item Транспортними засобами Університету (за погодженням з адміністрацією).
        \end{enumerate}

    8.4.3. ЦВКс забезпечує збереження та раціональне використання наданого майна.

    8.4.4. Передача майна ЦВКс здійснюється на підставі відповідних актів з визначенням відповідальної особи.

\subsection*{8.5. Матеріально-технічне забезпечення виборів}
\addcontentsline{toc}{subsection}{8.5. Матеріально-технічне забезпечення виборів}
    8.5.1. ЦВКс забезпечує виготовлення та розповсюдження виборчих матеріалів:

        \begin{enumerate}[label=\alph*)]
            \item Виборчих бюлетенів (для паперового голосування);
            \item Протоколів засідань та голосування;
            \item Інформаційних матеріалів для виборців;
            \item Посвідчень для членів виборчих комісій та спостерігачів.
        \end{enumerate}

    8.5.2. Для електронного голосування ЦВКс забезпечує:

        \begin{enumerate}[label=\alph*)]
            \item Функціонування електронних систем голосування;
            \item Навчання операторів електронних систем;
            \item Технічну підтримку під час голосування;
            \item Резервне копіювання електронних даних.
        \end{enumerate}

    8.5.3. Всі виборчі матеріали виготовляються за рахунок коштів, передбачених у кошторисі ЦВКс.

\subsection*{8.6. Звітність та прозорість}
\addcontentsline{toc}{subsection}{8.6. Звітність та прозорість}
    8.6.1. ЦВКс забезпечує повну прозорість використання коштів та майна шляхом:

        \begin{enumerate}[label=\alph*)]
            \item Оприлюднення кошторису та звітів про витрати;
            \item Надання доступу до фінансової документації зацікавленим особам;
            \item Проведення відкритих тендерів при закупівлі товарів та послуг;
            \item Регулярного інформування студентської спільноти про використання ресурсів.
        \end{enumerate}

    8.6.2. Порядок оприлюднення фінансової звітності ЦВКс визначається рішеннями КСУ та внутрішніми процедурами органів студентського самоврядування.